\documentclass[a4paper,12pt]{article}

\usepackage{polyglossia}
\setmainlanguage{russian}
\usepackage{fontspec}
\setmainfont{Times New Roman}
\usepackage{geometry}
\usepackage{setspace}
\usepackage{amsmath}
\usepackage{amssymb}
\usepackage{graphicx} 
\usepackage{listings}
\usepackage{xcolor}
\setmonofont{Courier New}
\newfontfamily{\cyrillicfonttt}{Courier New}


\geometry{top=20mm, bottom=20mm, left=25mm, right=15mm}

\lstset{
    language=Python,
    basicstyle=\ttfamily\scriptsize,
    keywordstyle=\color{blue},
    stringstyle=\color{red},
    commentstyle=\color{green!60!black},
    numbers=left,
    numberstyle=\tiny\color{gray},
    breaklines=true,
    frame=single
}

\begin{document}

\noindent
\begin{minipage}[c]{0.22\textwidth}
    \centering
    
\includegraphics[width=\textwidth]{../bmstu_logo.jpg} 
\end{minipage}%
\hfill
\begin{minipage}[c]{0.75\textwidth}
    \centering
    \setstretch{0.95}
    {\small\bfseries Министерство науки и высшего образования Российской Федерации} \\[2pt]
    {\small\bfseries Федеральное государственное автономное образовательное учреждение} \\[2pt]
    {\small\bfseries высшего образования} \\[2pt]
    {\small\bfseries «Московский государственный технический университет} \\[2pt]
    {\small\bfseries имени Н.Э. Баумана} \\[2pt]
    {\small\bfseries (национальный исследовательский университет)»} \\[2pt]
    {\small\bfseries (МГТУ им. Н.Э. Баумана)}
\end{minipage}

\vspace{12pt}
\hrule height 1pt
\vspace{2pt}
\hrule height 2pt

\vspace{10mm}

\noindent
\setstretch{1.2}
\textbf{ФАКУЛЬТЕТ} \underline{\makebox[138mm][l]{\hspace{15mm}\textbf{ИНФОРМАТИКА И СИСТЕМЫ УПРАВЛЕНИЯ}}} \\[10pt]
\textbf{КАФЕДРА} \underline{\makebox[143mm][l]{\hspace{12mm}\textbf{СИСТЕМЫ ОБРАБОТКИ ИНФОРМАЦИИ И УПРАВЛЕНИЯ}}}

\vspace{25mm}

\begin{center}
    \setstretch{1.3}
    {\Large Отчет по лабораторной работе № 6} \\[12pt]
    {\Large\bfseries «Необходимые и достаточные условия} \\[4pt]
    {\Large\bfseries существования условного экстремума»} \\[20mm]
    {\large по дисциплине «Технологии машинного обучения»}
\end{center}

\vspace{35mm}

\noindent
\setstretch{1.1}
\begin{tabular}{@{}llcc}
    \large Студент & \underline{\makebox[28mm][c]{\textbf{ИУ5-64Б}}} & \underline{\makebox[55mm]{}} & \underline{\makebox[50mm][c]{\textbf{Н.А. Чернев}}} \\[2pt]
    & \multicolumn{1}{c}{\scriptsize (Группа)} & \scriptsize (Подпись, дата) & \scriptsize (И.О. Фамилия) \\[20pt]
    \large Преподаватель & & \underline{\makebox[55mm]{}} & \underline{\makebox[50mm][c]{\textbf{Ю.Е. Гапанюк}}} \\[2pt]
    & & \scriptsize (Подпись, дата) & \scriptsize (И.О. Фамилия)
\end{tabular}

\vfill

\begin{center}
    \large 2026 г.
\end{center}

\newpage

\pagestyle{plain}

\section*{Цель работы}
\begin {enumerate}
\item Исследование необходимых и достаточных условий существования экстремума функции с учетом ограничений (условный экстремум).
\item Вычисление экстремумов функции.
\end {enumerate}
\section*{Постановка задачи}
Исследование необходимых и достаточных условий существования экстремума функции с учетом ограничений (условный экстремум) с целевой функцией:
$$f(x) = (10\ln x_1 - 2)^2 + 10(x_2^2 + 2)^2 - 10 \rightarrow \text{min}$$
при ограничениях:
$$\begin{cases}
x_1^2 + x_2^2 \le 2 \\
x_1 > 0, \quad x_2 \ge 0
\end{cases}$$

\section*{Аналитическое решение}

\subsection*{Условия Куна-Таккера}
Функция Лагранжа ($\lambda_0=1$):
$$L(x, \lambda) = (10\ln x_1 - 2)^2 + 10(x_2^2 + 2)^2 - 10 + \lambda_1(x_1^2 + x_2^2 - 2)$$

Уравнения стационарности:
$$\frac{\partial L}{\partial x_1} = \frac{20(10\ln x_1 - 2)}{x_1} + 2\lambda_1 x_1 = 0$$
$$\frac{\partial L}{\partial x_2} = 40x_2(x_2^2 + 2) + 2\lambda_1 x_2 = 0$$

\subsection*{Стационарная точка}
Рассмотрев случай активного ограничения $g_3(x)$ ($x_2 = 0$) и пассивного $g_1(x)$ ($\lambda_1 = 0$):

Из второго уравнения: $\lambda_3 = 0$.

Из первого: $10\ln x_1 - 2 = 0 \implies x_1 = e^{0.2} \approx 1.221$

Проверка допустимости: $(e^{0.2})^2 + 0^2 - 2 = e^{0.4} - 2 \approx -0.508 < 0$ ✓

\textbf{Решение:} $x^* = (e^{0.2}, 0) \approx (1.221, 0)$ при $f(x^*) = 30$

\section*{Численное решение методами Python}

\begin{lstlisting}
import numpy as np
from scipy.optimize import minimize
import matplotlib.pyplot as plt

def objective(x):
    x1 = max(x[0], 1e-9)
    return (10 * np.log(x1) - 2)**2 + 10 * (x[1]**2 + 2)**2 - 10

def constraint1(x):
    return 2 - (x[0]**2 + x[1]**2)

bounds = ((1e-9, None), (0, None))
cons = {'type': 'ineq', 'fun': constraint1}

# Численная минимизация
res = minimize(objective, [1.0, 0.5], method='SLSQP', 
               bounds=bounds, constraints=cons)

print(f"x*: [{res.x[0]:.4f}, {res.x[1]:.4f}]")
print(f"f(x*): {res.fun:.4f}")
\end{lstlisting}

\section*{Визуализация}

\begin{lstlisting}
# Построение контурных линий
x1_vals = np.linspace(0.01, 2.0, 400)
x2_vals = np.linspace(-0.5, 2.0, 400)
X1, X2 = np.meshgrid(x1_vals, x2_vals)

Z = (10 * np.log(X1) - 2)**2 + 10 * (X2**2 + 2)**2 - 10

plt.figure(figsize=(10, 8))
contours = plt.contour(X1, X2, Z, levels=[25, 30, 35, 50, 100, 200])
plt.clabel(contours, inline=True, fontsize=8)

# Допустимое множество
theta = np.linspace(0, np.pi/2, 200)
x1_c = np.sqrt(2) * np.cos(theta)
x2_c = np.sqrt(2) * np.sin(theta)

plt.fill_between(x1_c, x2_c, color='green', alpha=0.15, label='Допустимое множество')
plt.plot(x1_c, x2_c, 'g--', linewidth=2)
plt.plot(res.x[0], res.x[1], 'ro', markersize=10, label=f'Минимум x*')
plt.axhline(0, color='black', linewidth=0.8)
plt.axvline(0, color='black', linewidth=0.8)

plt.title('Линии уровня и допустимое множество (Вариант 10)')
plt.xlabel('$x_1$')
plt.ylabel('$x_2$')
plt.xlim(0, 1.8)
plt.ylim(-0.2, 1.6)
plt.grid(True, alpha=0.3)
plt.legend()
plt.savefig('fig_1.png', dpi=150, bbox_inches='tight')
plt.show()
\end{lstlisting}

\begin{figure}[h!]
    \centering
    \includegraphics[width=0.85\textwidth]{fig_1.png}
    \newline
    \small Рисунок 1 --- Линии уровня и допустимое множество для задачи условной оптимизации (Вариант 10)
    \newline
\end{figure}

\section*{Результаты}

\begin{itemize}
    \item \textbf{Аналитическое решение:} $x^* = (1.2214, 0)$, $f(x^*) = 30$
    \item \textbf{Численное решение (SLSQP):} совпадает с аналитическим
    \item \textbf{Гессиан в $x^*$:} положительно определен (локальный минимум)
\end{itemize}

\section*{Выводы}

\begin{enumerate}
    \item Метод SLSQP надежно находит локальный экстремум при наличии ограничений
    \item Активное ограничение $x_2 = 0$ существенно для оптимального решения
    \item Нарушение граничного условия $x_1 > 0$ приводит к недопустимым решениям
    \item Графическое представление помогает верифицировать результаты оптимизации
\end{enumerate}

\end{document}
